\exercisesection{不变子空间}

\begin{enumerate}
\def\labelenumi{\arabic{enumi})}
\tightlist
\item
  设 \((a_{n})_{n\in \mathbf{N}}\) 为一列复数，满足：存在实数 \(A\geqslant 0\)，使得 \(n|a_n|\leqslant A\)（对所有 \(n\in \mathbf{N}\)）。对 \(z\in \mathbf{C}\)（模 \(<  1\)），置 \(\varphi (z) = \sum_{n\geqslant 0}a_nz^n\)。
\end{enumerate}

我们假设存在 \(\ell \in \mathbf{C}\) 使得

\[
\lim_{\substack{x\in \mathbf{R}_{+}\\ x <   1}}\varphi (x) = \ell .
\]

对 \(t\in \mathbf{R}_+^*\)，置

\[
g (t) = \sum_ {0 \leqslant n \leqslant t ^ {- 1}} a _ {n}, \quad \mathrm{F} (t) = \varphi (e ^ {- t}).
\]

我们用 \(dx\) 表示 \(\mathbf{R}\) 上的 Lebesgue 测度，以及它在 \(\mathbf{R}_{+}^{*}\) 上的限制。

\begin{enumerate}
\def\labelenumi{\alph{enumi})}
\tightlist
\item
  设 \(s \leqslant t\) 为 \(\mathbf{R}_{+}^{*}\) 中的元素。我们有
\end{enumerate}

\[
| g (s) - g (t) | \leqslant | a _ {[ t ^ {- 1} + 1 ]} | + A \log \left(\frac {t}{s}\right).
\]

\begin{enumerate}
\def\labelenumi{\alph{enumi})}
\setcounter{enumi}{1}
\item
  函数 g 当 \(t \rightarrow 0\) 时是缓变的。
\item
  对每个 \(t \in R_{+}^{*}\)，函数 \(x \mapsto g(t/x)xe^{-x}\) 在 \(R_{+}^{*}\) 上关于 Haar 测度 \(x^{-1}dx\) 是可积的。
\item
  对 \(0 < x < 1\)，有
\end{enumerate}

\[
\varphi (x) = \sum_ {n \geqslant 0} g \left(\frac {1}{n}\right) \left(x ^ {n} - x ^ {n + 1}\right).
\]

\begin{enumerate}
\def\labelenumi{\alph{enumi})}
\setcounter{enumi}{4}
\tightlist
\item
  设 \(f(x) = xe^{-x}\)（对 \(x \in R_{+}^{*}\)）。我们有 \(f \in \mathrm{L}^{1}(\mathbf{R}_{+}^{*}, dx/x)\)；它的积分等于 1，并且我们有 F = g * f。
\end{enumerate}

\(f)\) 函数 \(g\) 属于 \(\mathrm{L}^{\infty}(\mathbf{R}_{+}^{*})\)。（证明 \(\mathrm{F}-g\) 是有界的（利用 \(a\) ) 以及假设）。）

\(g)\) 函数 \(f\) 在群 \(\mathbf{R}_{+}^{*}\) 上的 Fourier 变换不取零值。（把 \(\mathcal{F}(f)\) 等同于函数 \(y\mapsto\Gamma(1+iy)\) 并利用 FVR, VII。）

\begin{enumerate}
\def\labelenumi{\alph{enumi})}
\setcounter{enumi}{7}
\tightlist
\item
  我们有
\end{enumerate}

\[
\sum_ {n = 0} ^ {+ \infty} a _ {n} = \ell .
\]

这个命题就是 Littlewood 的“Tauber 型”定理。关于序列 \((a_n)\) 的假设，即 \(a_n = \mathrm{O}(1/n)\)，还可以减弱：只需存在常数 \(A \geqslant 0\)，使得 \(na_n \geqslant -A\)（对每个整数 \(n \geqslant 1\)）（“Hardy–Littlewood Tauber 型定理”，例如参看 J. KOREVAAR, Tauberian theory, Grundlehren math. wiss. 329, Springer (2004), Th. 7.4, 第 16 页；读者应注意不要把这个定理与 FVR, I, 第 50 页, 练习 18 中的定理混淆）。

\begin{enumerate}
\def\labelenumi{\arabic{enumi})}
\setcounter{enumi}{1}
\item
  设 F 为 \(\widehat{\mathbf{G}}\) 的闭子集。设 K 为 \(\widehat{\mathbf{G}}\) 的紧子集，使得 \(\mathrm{F} \cap \mathrm{K} = \varnothing\)。证明存在 G 上可积的连续函数 \(g\)，使得 \(\mathcal{F}(g)\) 在 F 上等于 0，在 K 上等于 1。（设 \(f \in \mathrm{L}^{1}(\mathrm{G})\) 满足 \(\mathcal{F}(f)\) 在 F 上为零且在 K 上等于 1；考察 \(f * \varphi\)，其中 \(\varphi\) 是形如 \(h * h'\) 的函数的 Fourier 变换（\(h \in \mathcal{K}(\widehat{\mathrm{G}})\)，\(h' \in \mathcal{K}(\widehat{\mathrm{G}})\)，且 \(h * h' = 1\) 在 K 上）。）
\item
  设 \(g\) 是定义在 \(]0, +\infty[\) 上的复值函数，关于 Lebesgue 测度可积。我们假设 \(\int_0^{+\infty} g(t)dt = 1\)，且 \(\int_0^{+\infty} g(t)t^{ix}dt \neq 0\)（对所有 \(x \in \mathbf{R}\)）。设 \(f\) 是 \(]0, +\infty[\) 上的复值可测有界函数，使得 \(\frac{1}{x}\int_0^{+\infty} g(t/x)f(t)dt\) 趋于一个有限极限 \(\ell\)（当 \(x\) 趋于 \(+\infty\)）。
\end{enumerate}

\begin{enumerate}
\def\labelenumi{\alph{enumi})}
\tightlist
\item
  对定义在 \(]0, +\infty[\) 上的、关于 Lebesgue 测度可积且积分为 1 的任意复值函数 h，我们有
\end{enumerate}

\[
\lim _ {x \to + \infty} \frac {1}{x} \int_ {0} ^ {+ \infty} h \left(\frac {t}{x}\right) f (t) d t = \ell .
\]

特别地，

\[
\lim _ {x \rightarrow + \infty} \frac {1}{x} \int_ {0} ^ {x} f (t) d t = \ell .
\]

\((\text{设 } g_1(t) = tg(t), h_1(t) = th(t), \text{ 我们有 } g_1 \in \mathrm{L}^1(\mathbf{R}_+^*) ; \text{ Fourier 变换 } \mathcal{F}(g_1) \text{ 不取零值且 } \check{g}_1 * f \text{ 趋于 } \ell, \text{ 因此 } h_1 * f \text{ 趋于 } \ell.)\) b) 若 \(f(t)\) 在 \(\mathbf{R}_+^*\) 上是缓变的（当 \(t\) 趋于 \(+\infty\)），则 \(f\) 趋于 \(\ell\)（当 \(t\) 趋于 \(+\infty\)）。

\begin{enumerate}
\def\labelenumi{\arabic{enumi})}
\setcounter{enumi}{3}
\tightlist
\item
  设 G = R。对每个实数 \(\alpha > 1\)，我们用下式定义 \(f_{\alpha} \in L^{1}(\mathbf{R})\)
\end{enumerate}

\[
f _ {\alpha} (x) = \left\{ \begin{array}{l l} 2 & \text { if } 0 <   x <   1 \\ 1 & \text { if } 1 \leqslant x <   \alpha \\ 0 & \text { otherwise. } \end{array} \right.
\]

\begin{enumerate}
\def\labelenumi{\alph{enumi})}
\item
  函数 \(f * \varepsilon_{x}\)（\(x \in R\)）在 \(\mathrm{L}^{1}(\mathbf{R})\) 中构成完全集，当且仅当 \(\alpha\) 是无理数。
\item
  函数 \(f * \varepsilon_x\)（\(x \in \mathbf{R}\)）在 \(L^2(\mathbf{R})\) 中构成完全集。
\end{enumerate}

\begin{enumerate}
\def\labelenumi{\arabic{enumi})}
\setcounter{enumi}{4}
\tightlist
\item
  设 \(f \in \mathrm{L}^{\infty}([0, +\infty([)\text{ 且 } k \in \mathbf{R}_{+}^{*}. \text{ 若}\)
\end{enumerate}

\[
\lim _ {x \to + \infty} \frac {1}{x} \int_ {0} ^ {+ \infty} e ^ {- t / x} f (t) d t = \ell
\]

则

\[
\lim _ {x \rightarrow + \infty} \frac {k}{x ^ {k}} \int_ {0} ^ {x} (x - t) ^ {k - 1} f (t) d t = \ell
\]

（利用函数 \(g(t) = e^{-t}\)，以及 \(h(t) = k(1 - t)^{k-1}\)（对 \(0 \leqslant t \leqslant 1\)；\(h(t) = 0\) 对 \(t > 1\)）。）

\begin{enumerate}
\def\labelenumi{\arabic{enumi})}
\setcounter{enumi}{5}
\tightlist
\item
  设 \(f \colon \mathbf{R} \to \mathbf{R}\) 与 \(g: \mathbf{R} \to \mathbf{C}\) 是两个函数。我们假设 \(f \in \mathrm{L}^{\infty}(\mathbf{R})\) 且 \(g \in \mathrm{L}^{1}(\mathbf{R})\)，\(\int_{\mathbf{R}} g(x) dx = 1\)，并且 \(\mathcal{F}(g)\) 不取零值。我们假设 \(f\) 是缓减的，即：对满足 \(y > x\)（\(y - x\) 趋于 0 且 \(x \to +\infty\)）的 y、x，有
\end{enumerate}

\[
\liminf (f (y) - f (x)) \geqslant 0.
\]

若 \((g*f)(x)\) 趋于 \(\ell\)（当 x 趋于 \(+\infty\)），则 \(f(x)\) 趋于 \(\ell\)（当 x 趋于 \(+\infty\)）。

\begin{enumerate}
\def\labelenumi{\arabic{enumi})}
\setcounter{enumi}{6}
\tightlist
\item
  设 \(g: R \rightarrow C\) 为连续函数，满足：
\end{enumerate}

\[
\sum_ {n = - \infty} ^ {+ \infty} \sup _ {n \leqslant t \leqslant n + 1} | g (t) | <   + \infty .
\]

\begin{enumerate}
\def\labelenumi{\alph{enumi})}
\item
  我们有 \(g \in \mathrm{L}^{1}(\mathbf{R})\)。
\item
  设 g 的积分为 1 且 \(\mathcal{F}(g)\) 不取零值。设 \(\mu\) 为 R 上的复测度，使得 \(|\mu|([x,x+1])\) 当 x 取遍 R 时是有界的。若 \((g*\mu)(x)\) 趋于 \(\ell\)（当 x 趋于 \(+\infty\)），则对每个满足下述条件的连续函数 \(h:R\to C\)
\end{enumerate}

\[
\sum_ {n = - \infty} ^ {+ \infty} \sup _ {n \leqslant t \leqslant n + 1} | h (t) | <   + \infty ,
\]

（积分为 1），\((h * \mu)(x)\) 当 \(x\) 趋于 \(+\infty\) 时的极限等于 \(\ell\)。（证明 \(h * \mu\) 是缓变的，并且 \((g * (h * \mu))(x)\) 趋于 \(\ell\)（当 \(x\) 趋于 \(+\infty\)）。）

\begin{enumerate}
\def\labelenumi{\arabic{enumi})}
\setcounter{enumi}{7}
\tightlist
\item
  设 \(g\) 是满足练习 3 的假设的函数。此外，设 \(g \geqslant 0\) 且 \(\liminf_{x \to 0} g(x) > 0\)。设 \(f\) 是 \(]0, +\infty[\) 上的有下界的实值可测函数。若
\end{enumerate}

\[
\lim _ {x \to + \infty} \frac {1}{x} \int_ {0} ^ {+ \infty} g \left(\frac {t}{x}\right) f (t) d t = \ell
\]

则

\[
\sigma (x) = \frac {1}{x} \int_ {0} ^ {x} f (t) d t
\]

趋于 \(\ell\)（当 \(x\) 趋于 \(+\infty\)）。（证明 \(\sigma\) 对 \(x \geqslant 1\) 是有界的，再证明 \(\sigma\) 是缓减的，然后证明

\[
\frac {1}{x} \int_ {0} ^ {+ \infty} g \left(\frac {t}{x}\right) \sigma (t) d t
\]

趋于 \(\ell\)（当 \(x\) 趋于 \(+\infty\)）。）

\begin{enumerate}
\def\labelenumi{\arabic{enumi})}
\setcounter{enumi}{8}
\tightlist
\item
  设 \(\varphi\colon[0,+\infty[\to\mathbf{R}\) 是非负的递增函数，使得函数 \(t\mapsto e^{-\sigma t}\varphi(t)\) 当 \(\sigma>1\) 时关于 Lebesgue 测度可积。对每个 \(s\in\mathbf{C}\)（满足 \(\mathcal{R}(s)>1\)），置
\end{enumerate}

\[
f (s) = \int_ {0} ^ {+ \infty} e ^ {- s t} \varphi (t) d t
\]

（“\(\varphi\) 的 Laplace 变换”）。

假设存在 \(A \in C\) 具有下列性质：当 \(\sigma\) 通过 \textgreater{} 1 的值趋于 1 时，函数

\[
t \mapsto f (\sigma + i \tau) - \frac {\mathrm{A}}{\sigma + i \tau - 1} \qquad \qquad \tau \in \mathbf {R}
\]

在 R 的每个紧子集上一致收敛于一个函数 g。则

\[
\lim _ {t \to + \infty} \varphi (t) e ^ {- t} = A
\]

（“Ikehara Tauber 型定理”）。

（我们置 \(a(t) = e^{-t}\varphi(t)\)（对 \(t > 0\)），\(a(t) = 0\)（对 \(t \leqslant 0\)），\(A(t) = A\)（对 \(t > 0\)），\(A(t) = 0\)（对 \(t \leqslant 0\)）。设 \(\lambda > 0\)。我们置：

\[
k _ {\lambda} (t) = \lambda \sqrt {\frac {2}{\pi}} \Big (\frac {\sin 2 \lambda t}{2 \lambda t} \Big) ^ {2}
\]

\[
\mathrm{K} _ {\lambda} (t) = \left\{ \begin{array}{l l} 1 - \left| \frac {t}{2 \lambda} \right| & \text { if } | t | \leqslant 2 \lambda \\ 0 & \text { if } | t | > 2 \lambda . \end{array} \right.
\]

利用 II，第 262 页，练习 1，Lebesgue–Fubini 定理以及 Plancherel 定理，证明：

\[
\lim _ {\varepsilon \to 0} \frac {1}{\sqrt {2 \pi}} \int_ {- \infty} ^ {+ \infty} k _ {\lambda} (x - t) (a (t) - \mathrm{A} (t)) e ^ {- \varepsilon t} d t = \frac {1}{\sqrt {2 \pi}} \int_ {- 2 \lambda} ^ {2 \lambda} \mathrm{K} _ {\lambda} (y) e ^ {- i x y} g (y) d y
\]

并推出

\[
\lim _ {x \to + \infty} \frac {1}{\sqrt {2 \pi}} \int_ {- \infty} ^ {+ \infty} k _ {\lambda} (x - t) a (t) d t = \mathrm{A}.
\]

然后证明 \(a\) 是有界的且缓减的。最后可以应用练习 6。）

\begin{enumerate}
\def\labelenumi{\arabic{enumi})}
\setcounter{enumi}{9}
\tightlist
\item
  我们在 \(\mathbf{N}^{*}\) 上定义 von Mangoldt 函数如下：\(\Lambda(n)=0\)（若 n 不是形如 \(p^{k}\) 的数，其中 p 是素数且 \(k\geqslant1\)），而 \(\Lambda(p^{k})=\log(p)\)。
\end{enumerate}

我们用 \(\zeta\) 表示 \(\mathbf{C} - \{1\}\) 上的全纯函数，它到满足如下条件的 \(s \in \mathbf{C}\) （\(\mathcal{R}(s) > 1\)）组成的集合上的限制与 Riemann zeta 函数重合（II，第 266 页，练习 16）。

对 \(x > 0\)，我们置

\[
\psi (x) = \sum_ {1 \leqslant n \leqslant x} \Lambda (n).
\]

\begin{enumerate}
\def\labelenumi{\alph{enumi})}
\tightlist
\item
  证明：对所有 \(s \in \mathbf{C}\)（满足 \(\mathcal{R}(s) > 1\)），有
\end{enumerate}

\[
\sum_ {n \geqslant 1} \Lambda (n) n ^ {- s} = - \frac {\zeta^ {\prime} (s)}{\zeta (s)}.
\]

\begin{enumerate}
\def\labelenumi{\alph{enumi})}
\setcounter{enumi}{1}
\tightlist
\item
  我们有
\end{enumerate}

\[
\lim _ {x \to + \infty} \frac {\psi (x)}{x} = 1.
\]

（利用 II，第 266 页，练习 16 以及练习 9）。

\begin{enumerate}
\def\labelenumi{\alph{enumi})}
\setcounter{enumi}{2}
\tightlist
\item
  设 \(\pi(x)\) 为 \(\leqslant x\) 的素数的个数。证明
\end{enumerate}

\[
\pi (x) \sim \frac {x}{\log (x)}
\]

（当 \(x \to +\infty\)）（“素数定理”，归于 Hadamard 与 de la Vallée Poussin）。

\begin{enumerate}
\def\labelenumi{\arabic{enumi})}
\setcounter{enumi}{10}
\tightlist
\item
  我们如下定义 Möbius 函数 \(\mu\)（在 \(\mathbf{N}^*\) 上）：\(\mu(n) = (-1)^k\)（若 \(n\) 是 \(k\) 个两两不同的素因子之积），否则 \(\mu(n) = 0\)。
\end{enumerate}

对 \(x > 0\)，我们置

\[
\mathrm{M} (x) = \sum_ {n \leqslant x} \mu (n), \quad f (x) = \frac {\mathrm{M} (x)}{x}.
\]

\begin{enumerate}
\def\labelenumi{\alph{enumi})}
\tightlist
\item
  证明：对 \(\mathcal{R}(s) > 1\)，有
\end{enumerate}

\[
\sum_ {n \geqslant 1} \mu (n) n ^ {- s} = \frac {1}{\zeta (s)}.
\]

\begin{enumerate}
\def\labelenumi{\alph{enumi})}
\setcounter{enumi}{1}
\item
  证明 f 是有界的，且 \(f(y) - f(x) = \mathrm{O}((y - x)/x)\)（对 \(y \geqslant x \geqslant 1\)），从而 f 在群 \(R_{+}^{*}\) 上当 x 趋于 \(+\infty\) 时是缓变的。
\item
  设 \(a, b \in R_{+}^{*}\)。对 \(x \textgreater{} 0\)，设 {[}x{]} 为 x 的整数部分。我们置 \(g_{0}(t) = [t^{-1}]\)（对 \(t \in R_{+}^{*}\)），并置
\end{enumerate}

\[
g (t) = 2 g _ {0} (t) - a g _ {0} (a t) - b g _ {0} (b t).
\]

函数 g 在 \(R_{+}^{*}\) 上关于 Lebesgue 测度可积。若 \(s \in C\) 且 \(\mathcal{R}(s) > 0\)，则

\[
\int_ {0} ^ {+ \infty} g (t) t ^ {s} d t = (2 - a ^ {- s} - b ^ {- s}) \frac {\zeta (1 + s)}{1 + s}.
\]

\(d)\) 我们有 \(\int_{0}^{+\infty} g(t)t^{ix} dt \neq 0\)（对所有 \(x \in \mathbf{R}\)）。（利用 II，第 266 页，练习 16 的 \(j\) )。）

\begin{enumerate}
\def\labelenumi{\alph{enumi})}
\setcounter{enumi}{4}
\tightlist
\item
  证明 \(\int_0^{+\infty} g(t / x)f(t)dt = o(x)\)（当 \(x \to +\infty\)）。推出 \(\mathrm{M}(x) = o(x)\)（当 \(x \to +\infty\)）。
\end{enumerate}

（可以初等地证明这个结果与素数定理等价，参看练习 10；例如参看 H. IWANIEC and E. KOWALSKI, Analytic number theory, A.M.S Coll. Publ. 53 (2004), 第 31--32 页。）

\begin{enumerate}
\def\labelenumi{\arabic{enumi})}
\setcounter{enumi}{11}
\item
  设 J 为满足如下条件的 \(f \in \mathcal{S}(\mathbf{R}^3)\) 组成的集合：\(\mathcal{F}(f)\) 在球面 \(\mathbf{S}_2\) 上为零。设 I 为满足如下条件的 \(f \in \mathrm{J}\) 组成的集合：\((\partial/\partial y_1)(\mathcal{F}(f))\) 在 \(\mathbf{S}_2\) 上为零。设 \(\overline{\mathrm{I}}\) 与 \(\overline{\mathrm{J}}\) 为 I 与 J 在 \(\mathrm{L}^1(\mathbf{R}^3)\) 中的闭包。则我们有 \(\mathrm{V}(\overline{\mathrm{I}}) = \mathrm{V}(\overline{\mathrm{J}}) = \mathbf{S}_2\)，但 \(\overline{\mathrm{I}} \neq \overline{\mathrm{J}}\)。（设 \(\mu\) 为 \(\mathbf{S}_2\) 上质量为 1、在 \(\mathbf{R}^3\) 的正交群下不变的正测度；利用 II，第 269 页，练习 20 的 c)，证明 \(f(t) = t_1\mathcal{F}(\mu)(t)\) 是 \(\mathrm{L}^\infty(\mathbf{R}^3)\) 中与 I 正交但不与 J 正交的元素。）
\item
  我们称 \(\widehat{\mathbf{G}}\) 中的 C 集是指 \(\widehat{\mathbf{G}}\) 的具有下列性质的闭子集 E：若 \(f\in \mathrm{L}^{1}(\mathrm{G})\)，若 \(\mathcal{F}(f)\) 在 E 上为零，且若 \(\varepsilon >0\)，则存在函数 \(g\in \mathrm{L}^{1}(\mathrm{G})\)，使得 \(\| f - f*g\| \leqslant \varepsilon\) 且 \(\operatorname {Supp}(\mathcal{F}(g))\) 是紧的并与 E 不相交。
\end{enumerate}

\begin{enumerate}
\def\labelenumi{\alph{enumi})}
\item
  若 E 是 \(\widehat{G}\) 中的 C 集，则恰存在 \(\mathrm{L}^1 (\mathrm{G})\) 的唯一闭理想 I 使得 \(\mathrm{V(I)} = \mathrm{E}\)。
\item
  \(\widehat{\mathrm{G}}\) 的每个单点子集都是 C 集。
\item
  两个 C 集的并是 C 集。
\item
  由 C 集经平移得到的每个集合都是 C 集。
\item
  设 H 为 \(\widehat{\mathbf{G}}\) 的闭子群，E 为 H 的闭子集，E' 为 E 相对于 H 的边界。若 E' 是相对于 \(\widehat{\mathbf{G}}\) 的 C 集，则 E 是关于 \(\widehat{\mathbf{G}}\) 的 C 集。（利用 II，第 270 页，练习 21。）
\item
  在 \(R^{n}\) 中，有限个闭半空间的交 E 是 C 集。（对 E 生成的仿射子空间的维数用归纳法，利用 b)、c)、d)、e)。
\end{enumerate}

\begin{enumerate}
\def\labelenumi{\arabic{enumi})}
\setcounter{enumi}{13}
\tightlist
\item
  设 E 为 \(R^{n}\) 的闭子集。假设存在 E 的一个内点 p，使得每条过 p 的直线与 E 的边界至多交于两点。则恰存在 \(L^{1}(\mathbf{R}^{n})\) 的唯一闭理想 I 使得 \(\mathcal{F}(f)\) 在 E 上为零。（把 f 看作 f 经以 p 为中心的位似所得到的函数在 \(L^{1}(\mathbf{R}^{n})\) 中的极限。）
\end{enumerate}

¶ 15) 我们用 μ 表示 G 的一个 Haar 测度，用 B 表示 e 在 \(\widehat{G}\) 中的邻域滤子。

\begin{enumerate}
\def\labelenumi{\alph{enumi})}
\tightlist
\item
  设 U 为 e 在 G 中的一个可积的开邻域。存在函数 \(a \in \mathrm{L}^{1}(\mathrm{G})\)，\(a \geqslant 0\)，积分为 1，在 U 之外为零，使得 \(\mathcal{F}(a) \in \mathrm{L}^{1}(\widehat{\mathrm{G}})\)，并且
\end{enumerate}

\[
\int_ {\mathrm{G}} a (x) ^ {2} d x \leqslant \frac {2}{\mu (\mathrm{U})}.
\]

（设 V 为 e 的紧邻域，使得 \(V \subset U\)，\(\mu(U) < \sqrt{2}\mu(V)\)；设 W 为 e 的紧对称邻域，使得 \(VW^{2} \subset U\)；取 \(a = \lambda f * g\)，其中 \(\lambda\) 是一个 \textgreater{} 0 的常数，f、g 是集合 VW 与 W 的特征函数。）

\begin{enumerate}
\def\labelenumi{\alph{enumi})}
\setcounter{enumi}{1}
\tightlist
\item
  若 \(f \in \mathrm{L}^{\infty}(\mathrm{G})\)，我们用 \(\mathrm{A}(f)\) 表示 \(\widehat{G}\) 中属于由 f 生成的、平移不变的弱闭向量子空间（\(\mathrm{L}^{\infty}(\mathrm{G})\) 的）的元素所成的集合。若 \(f, g \in \mathrm{L}^{\infty}(\mathrm{G})\)，则
\end{enumerate}

\[
\mathrm{A} (f g) \subset \overline {{\mathrm{A} (f) \mathrm{A} (g)}}.
\]

（设 U 为 \(e\) 的邻域。证明 \(fg\) 是 \(\widehat{G}\) 中属于 \(\mathrm{A}(f)\mathrm{UA}(g)\mathrm{U}\) 的元素的线性组合的弱极限。）

\begin{enumerate}
\def\labelenumi{\alph{enumi})}
\setcounter{enumi}{2}
\tightlist
\item
  设 \(g \in \mathrm{L}^{\infty}(\mathrm{G})\)，K 为 \(\widehat{G}\) 的紧子集，f 为 \(\mathrm{L}^{1}(\mathrm{G})\) 中的函数，使得 \(\mathcal{F}(f)\) 在 \(\mathrm{A} = \mathrm{A}(g)\) 上及在 K 之外为零。对每个 \(V \in B\)，设 \(\omega(V)\) 为 \(|\mathcal{F}(f)|\) 在 AV 上的上确界。若
\end{enumerate}

\[
\operatorname * {l i m i n f} _ {\mathfrak {B}} \frac {\omega (\mathrm{V}) ^ {2} \mu ((\mathrm{AV} - \mathrm{A}) \cap \mathrm{K})}{\mu (\mathrm{V})} = 0,
\]

则 \(\langle f,g\rangle=0\)。（设 \(\varepsilon>0\)。存在 G 的紧子集 H 使得 \(\int_{G-H}|f|dx\leqslant\varepsilon\)。于是存在 e 在 \(\widehat{G}\) 中的开邻域 U，使得 \(|\langle x,\widehat{x}\rangle-1|\leqslant\varepsilon\)（对 \(x\in H\)，\(\widehat{x}\in U\)）。把 a) 应用于 \(\widehat{G}\) 与 U，得到一个函数 \(a\in L^{1}(\widehat{G})\)。设 \(b=\mathcal{F}(a)\in L^{1}(G)\)。利用 b)，证明 \(\mathcal{F}(bg)\) 在 AU 之外为零。另一方面，

\[
\left| \int_ {\mathrm{G}} f g d x - \int_ {\mathrm{G}} f g b d x \right| \leqslant \varepsilon (\| f \| _ {1} + 2) \| g \| _ {\infty}
\]

并且，由于 \(f, gb \in \mathrm{L}^{2}(\mathrm{G})\)，

\[
\int_ {\mathrm{G}} f g b d x = \int_ {\widehat {\mathrm{G}}} \mathcal {F} (f) \mathcal {F} (g b) d \widehat {x} = \int_ {(\mathrm{AU-A}) \cap \mathrm{K}} \mathcal {F} (f) \mathcal {F} (g b) d \widehat {x}.
\]

把最后这个表达式的绝对值用下式界定：

\[
2 \| g \| _ {\infty} \frac {\omega (\mathrm{U}) \sqrt {\mu ((\mathrm{AU} - \mathrm{A}) \cap \mathrm{K})}}{\sqrt {\mu (\mathrm{U})}}
\]

即可完成证明。）

\(d)\) 设 I 为 \(\mathrm{L}^1 (\mathrm{G})\) 的闭理想，\(\mathrm{A} = \mathrm{V}(\mathrm{I})\)，\(f\) 为 \(\mathrm{L}^1 (\mathrm{G})\) 中的函数，使得 \(\mathcal{F}(f)\) 在 A 上为零，S 为 \(\mathcal{F}(f)\) 不为零的点所成的集合，\(\mathrm{S}'\) 为 S 的边界，B 为含于 \(\mathrm{S} \cap \mathrm{A}\) 的最大完满集。对 \(\mathrm{V} \in \mathfrak{B}\)，设 \(\omega (\mathrm{V})\) 为 \(|\mathcal{F}(f)|\) 在 AV 上的上确界。我们假设 B 的每个点在 \(\widehat{\mathrm{G}}\) 中有一个邻域 K，使得

\[
\operatorname * {l i m i n f} _ {\mathfrak {B}} \frac {\omega (\mathrm{V}) ^ {2} \mu ((\mathrm{AV-A}) \cap \mathrm{K})}{\mu (\mathrm{V})} = 0.
\]

则 \(f \in I\)。（利用 \(L^1(G)\) 满足 Ditkin 条件这一事实，并按 I，第 94 页，命题 5 的方式论证；设 \(G' = \widehat{G} \cup \{\infty\}\) 为 \(\widehat{G}\) 的 Alexandroff 紧化，并设 N 为满足如下条件的 \(\chi \in G'\) 所成的集合：\(\mathcal{F}(f)\) 不属于 \(\mathcal{F}(I)\)（在 \(\chi\) 的一个邻域中）。首先证明 \(N \subset B \cup \{\infty\}\)。然后，利用 \(c\) )，证明 \(N \subset \{\infty\}\)。最后，证明 \(N = \varnothing\)。）

\begin{enumerate}
\def\labelenumi{\alph{enumi})}
\setcounter{enumi}{4}
\tightlist
\item
  设 I 为 \(\mathrm{L}^{1}(\mathbf{R}^{n})\) 的闭理想，A = V(I)，f 为 \(\mathrm{L}^{1}(\mathbf{R}^{n})\) 中的函数，使得 \(\mathcal{F}(f)\) 在 A 上为零。对 h \textgreater{} 0，设 \(A_{h}\) 为 \(R^{n}\) 中在 A 之外且与 A 的距离小于 h 的点所成的集合，并设 \(\omega(h) = \sup_{x \in \mathrm{A}_{h}} |\mathcal{F}(f)(x)|\)。若
\end{enumerate}

\[
\liminf _ {h \to 0} \frac {\omega (h) ^ {2} \mu (\mathrm{A} _ {h})}{h ^ {n}} = 0
\]

则 \(f\in \mathrm{I}\)。（利用 \(d\) )。）

\(f)\) 设 \(f \in \mathrm{L}^{1}(\mathbf{R})\)，\(\alpha \in]0,1[\)，并设 \(\int_{\mathbf{R}} |y|^{\alpha}| f(y)| dy < +\infty\)。则存在常数 \(k\)，使得 \(|\mathcal{F}(f)(x+h)-\mathcal{F}(f)(x)| \leqslant kh^{\alpha}\)（对所有 \(x \in \mathbf{R}\) 以及每个 \(h \in \mathbf{R}\)）。

\(g)\) 设 I 为 \(\mathrm{L}^1 (\mathbf{R})\) 的闭理想，\(f\) 为 \(\mathrm{L}^1 (\mathbf{R})\) 中的函数，使得 \(\mathcal{F}(f)\) 在 \(\mathrm{V(I)}\) 上为零。若

\[
\int_ {\mathbf {R}} | y ^ {1 / 2} f (y) | d y <   + \infty ,
\]

则我们有 \(f \in I\)。（利用 e) 与 f)。）

\begin{enumerate}
\def\labelenumi{\arabic{enumi})}
\setcounter{enumi}{15}
\tightlist
\item
  设 \(f \in \mathrm{L}^{\infty}(\mathrm{G})\)。我们称 \(f\) 满足谱综合条件，如果它属于 \(\mathrm{Y}(\mathrm{A}(\mathrm{W}_f))\)，其中 \(\mathrm{W}_f\) 是 \(\mathrm{L}^{\infty}(\mathrm{G})\) 中包含 \(f\) 的所有平移不变弱闭子空间的交。
\end{enumerate}

\begin{enumerate}
\def\labelenumi{\alph{enumi})}
\item
  设 \(\mu \in \mathcal{M}^1 (\widehat{\mathrm{G}})\) 且 \(f = \overline{\mathcal{F}}_{\widehat{\mathrm{G}}}(\mu)\in \mathrm{L}^{\infty}(\mathrm{G})\)。则 \(\mathrm{A}(\mathrm{W}_f)\) 是 \(\mu\) 的支集。此外，\(f\) 满足谱综合条件。
\item
  若 G 是离散的，则每个函数 \(f \in \mathrm{L}^{2}(\mathrm{G})\) 都满足谱综合条件。
\end{enumerate}

\begin{enumerate}
\def\labelenumi{\arabic{enumi})}
\setcounter{enumi}{16}
\tightlist
\item
  对每个 \(p \geqslant 1\)，我们把 \(\mathrm{L}^p(\mathbf{R}_+)\) 与 \(\mathrm{L}^p(\mathbf{R})\) 的一个闭子空间等同。设 \(f \in \mathrm{L}^1(\mathbf{R}_+)\) 为由 \(f(x) = 2e^{-x}\)（对 \(x \in \mathbf{R}_+\)）定义的函数，并设 \(g \in \mathrm{L}^1(\mathbf{R})\) 为满足 \(g(x) = f(-x)\) 的函数。
\end{enumerate}

\begin{enumerate}
\def\labelenumi{\alph{enumi})}
\item
  证明 \(f + g = f * g\)。
\item
  设 A（相应地，B）为 \(\mathrm{L}^{1}(\mathbf{R})\) 中由 f（相应地，由 f 与 g）生成的子代数。证明 A 在 \(\mathrm{L}^{1}(\mathbf{R}_{+})\) 中稠密。（设 \(h \in \mathrm{L}^{\infty}(\mathbf{R}_{+})\) 满足：\(\langle h, f^{*n} \rangle = 0\)（对每个整数 \(n \geqslant 1\)），其中 \(f^{*n}\) 是 f 的 n 次迭次卷积；证明 Laplace 变换
\end{enumerate}

\[
\mathrm{F} (z) = \int_ {0} ^ {+ \infty} e ^ {- x z} h (x) d x,
\]

对所有 \(z \in C\)（满足 \(\mathcal{R}(z) > 0\)）都有定义的这个变换为零。）推出 B 在 \(\mathrm{L}^{1}(\mathbf{R})\) 中稠密。

\begin{enumerate}
\def\labelenumi{\alph{enumi})}
\setcounter{enumi}{2}
\item
  设 C 为 \(\mathrm{L}^{1}(\mathbf{R})\) 的包含 A 的真闭子代数。我们有 \(1 \in \mathrm{Sp}_{\mathrm{C}}(f)\)。（证明：否则，a) 以及 C 中的全纯函数演算将蕴含 \(g \in C\)。）
\item
  存在 C 的特征标 \(\chi\) 使得 \(\chi(f)=1\)；我们有
\end{enumerate}

\[
\chi (\varphi) = \int_ {0} ^ {+ \infty} e ^ {- x} \varphi (x) d x
\]

（对每个函数 \(\varphi\in\mathrm{L}^{1}(\mathbf{R}_{+})\)）。

\begin{enumerate}
\def\labelenumi{\alph{enumi})}
\setcounter{enumi}{4}
\item
  存在范数为 1 的测度 \(\mu \in \mathcal{M}^1 (\mathbf{R})\)，使得 \(\chi (\varphi) = \int \widehat{\varphi} d\mu\)（对每个函数 \(\varphi \in \mathrm{C}\)）。
\item
  \(\mu\) 的 Fourier 变换是函数 \(x \mapsto e^{-|x|}\)。
\item
  设 \(\varphi \in \mathrm{C}\)。对每个 \(y > 0\)，有
\end{enumerate}

\[
\int_ {\mathbf {R}} \varphi (x) (e ^ {- | x + y |} - e ^ {- | x | - y}) d x = 0.
\]

（利用关系式 \(\chi (\varphi *\psi) = \chi (\varphi)\chi (\psi)\)（对 \(\psi \in \mathrm{L}^1 (\mathbf{R}_+)\)）。）

\begin{enumerate}
\def\labelenumi{\alph{enumi})}
\setcounter{enumi}{7}
\tightlist
\item
  我们有 \(C = L^{1}(\mathbf{R}_{+})\)。（于是，子代数 \(L^{1}(\mathbf{R}_{+})\) 是 \(L^{1}(\mathbf{R})\) 的极大闭子代数。）
\end{enumerate}

\begin{enumerate}
\def\labelenumi{\arabic{enumi})}
\setcounter{enumi}{17}
\tightlist
\item
  设 G 为一个离散的有序交换拓扑群。我们把 G 的群运算写成加法，并用 0 表示它的单位元。我们用 \(G^{+}\) 表示 G 中满足 \(x \geqslant 0\) 的元素 x 所成的集合。
\end{enumerate}

\begin{enumerate}
\def\labelenumi{\alph{enumi})}
\item
  空间 \(L_{+}^{1}\)（由函数 \(f \in \mathrm{L}^{1}(\mathrm{G})\)（在 \(G^{+}\) 之外为零）组成）是 \(\mathrm{L}^{1}(\mathrm{G})\) 的闭子代数。
\item
  我们假设 G 不是 Archimedean 的（TG，V，第 16 页，§3 的练习 1）。代数 L \(_{+}^{1}\) 在 L \(^{1}\) (G) 中不是极大的。
\end{enumerate}

以下假设 G 是 Archimedean 的。设 A 为 \(\mathrm{L}^{1}(\mathrm{G})\) 的包含 \(L_{+}^{1}\) 的闭子代数。对每个 \(x \in G\)，设 \(\varphi_{x} \in \mathrm{L}^{1}(\mathrm{G})\) 为 \(\{x\}\) 的特征函数。它在 \(\mathrm{L}^{1}(\mathrm{G})\) 中可逆，其逆为 \(\varphi_{-x}\)。

\begin{enumerate}
\def\labelenumi{\alph{enumi})}
\setcounter{enumi}{2}
\tightlist
\item
  存在 \(y \in G_{+}\) 使得 \(\varphi_{y} \in A\) 且 \(\varphi_{-y} = \varphi_{y}^{-1} \notin A\)，并且存在 A 的特征标 \(\chi \neq 0\) 使得 \(\chi(y) = 0\)。
\end{enumerate}

\(d)\) 对 G 中每个 \(x > 0\)，我们有 \(\chi(x) = 0\)。（存在整数 \(k \geqslant 1\) 使得 \(-y + kx > 0\)。）

\begin{enumerate}
\def\labelenumi{\alph{enumi})}
\setcounter{enumi}{4}
\item
  存在范数为 1 的测度 \(\mu \in \mathcal{M}^1 (\mathrm{G})\)，使得 \(\chi\) 是 \(\mu\) 在代数 A（视为 \(\mathcal{C}(\mathrm{G})\) 的子空间）上的限制。（对 \(f\in \mathrm{A}\)，有 \(|\chi (f)|\leqslant \varrho (f)\)，而且 \(\varrho (f) = \| f\|_{\infty}\)。）
\item
  设 \(\widehat{\mu}\colon\mathrm{G}\to\mathbf{C}\) 为 \(\mu\) 的逆 Fourier 变换。我们有 \(\widehat{\mu}(x)=0\)（若 \(x>0\)）。
\item
  测度 \(\mu\) 是 G 上的计数测度。
\item
  子代数 \(L_{+}^{1}\) 在 \(L^{1}(G)\) 中是极大的。
\end{enumerate}

\begin{enumerate}
\def\labelenumi{\arabic{enumi})}
\setcounter{enumi}{18}
\tightlist
\item
  对 \(f \in \mathrm{L}^1(\mathbf{R})\)，我们用 \(\mathrm{A}_f\) 表示 \(\mathrm{L}^1(\mathbf{R})\) 中由 \(f\) 生成的闭子代数。
\end{enumerate}

\begin{enumerate}
\def\labelenumi{\alph{enumi})}
\tightlist
\item
  存在 \(f \in \mathrm{L}^{1}(\mathbf{R})\) 使得 \(A_{f}\) 是 \(\mathrm{L}^{1}(\mathbf{R})\) 的真极大子代数。
\end{enumerate}

在下文中，我们假设 \(f \in \mathrm{L}^{1}(\mathbf{R})\) 是使得 \(A_{f}\) 为 \(\mathrm{L}^{1}(\mathbf{R})\) 的真极大子代数的函数。

\begin{enumerate}
\def\labelenumi{\alph{enumi})}
\setcounter{enumi}{1}
\item
  我们记 \(S = \{0\} \cup \widehat{f}(\mathbf{R}) \subset \mathbf{C}\)。证明函数 \(z \mapsto \overline{z}\) 不是 S 上多项式函数的一致极限。（首先证明 C-S 不是连通的。）
\item
  \(\mathrm{L}^{1}(\mathbf{R})\) 中由 f 与 \(\widetilde{f}\) 生成的子代数在 \(\mathrm{L}^{1}(\mathbf{R})\) 中稠密。
\item
  f 的 Fourier 变换是单射的且在 R 上不取零值。（证明 S 不可能有两个有界的连通分支。）
\end{enumerate}

\backmatter
